\documentclass{article}
\usepackage{amsmath, amssymb}
% Allow the use of \cal as a shortcut for \mathcal
\let\cal\mathcal
\begin{document}

% we are trying to understand the standard model
% here are different renditions of it

% referential
\[
{\cal L}_{SM}
=
{\cal L}_{gauge}
+
{\cal L}_{fermions}
+
{\cal L}_{Yukawa}
+
{\cal L}_{Higgs}
\]

% expounded
\[
\mathcal{L} % lagrangian density
\quad = \quad 
-\frac{1}{4}F_{\mu\nu}^{\alpha}F^{\alpha\mu\nu} % kinetic term for gauge fields (electromagnetism/strong force)
\quad + \quad
i\overline{\psi}\gamma_{5}\psi + \text{h.c.} % fermion kinetic term (how matter particles move)
\quad + \quad
\overline{\psi_{L}}\gamma_{5}\psi_{R}\Phi + \text{h.c.} % yukawa coupling (gives particles mass via higgs)
\quad + \quad
|D_{\mu}\Phi|^{2} - V(\Phi) % higgs sector (spontaneous symmetry breaking)
\]

% full
\begin{align*}
\mathcal{L}_{\text{SM}}
\quad &= \quad 
-\frac{1}{4}G_{\mu\nu}^a G^{a\mu\nu} 
\quad - \quad
\frac{1}{4}W_{\mu\nu}^i W^{i\mu\nu}
\quad - \quad
\frac{1}{4}B_{\mu\nu}B^{\mu\nu} \\
&+ \qquad 
\overline{\psi}_L i\gamma^\mu D_\mu\psi_L 
\quad + \quad
\overline{\psi}_R i\gamma^\mu D_\mu\psi_R \\
&+ \qquad
- (y_{ij}\overline{\psi}_{L,i}\Phi\psi_{R,j} 
\quad + \quad
\text{h.c.}) \\
&+ \qquad
(D_\mu\Phi)^\dagger(D^\mu\Phi) 
\quad - \quad
\mu^2\Phi^\dagger\Phi 
\quad - \quad 
\lambda(\Phi^\dagger\Phi)^2 \\
&+ \qquad 
\text{gauge-fixing terms} 
\: + \:
\text{ghost terms} 
\: + \:
\theta\text{-terms} 
\: + \:
\text{neutrino mass terms (if included)}
\end{align*}



\end{document}